% Numeracja i wygląd zgodne z układem przekazanego szablonu.
\makeatletter
\def\numberline#1{\hb@xt@\@tempdima{#1.\hfil}}
\renewcommand*\@seccntformat[1]{\csname the#1\endcsname.\enspace}
\@addtoreset{equation}{section}
\renewcommand{\theequation}{\thesection.\arabic{equation}}
\@addtoreset{figure}{section}
\renewcommand{\thefigure}{\thesection.\arabic{figure}}
\@addtoreset{table}{section}
\renewcommand{\thetable}{\thesection.\arabic{table}}
\makeatother

\newenvironment{tabela}[3]{%
  \begin{table}[H]\centering\caption[#1]{#2}\medskip#3%
}{\end{table}}

\newcommand{\bezligaturkodu}{\ifdefined\pdfnoligatures\pdfnoligatures\font\fi}
\newcommand{\kod}[1]{\texttt{\bezligaturkodu\detokenize{#1}}}
\newcolumntype{L}[1]{>{\raggedright\arraybackslash}p{#1}}
\newcolumntype{Y}{>{\raggedright\arraybackslash}X}

% Listing pozostaje przy swoim opisie i nie jest dzielony pomiędzy strony.
\newcommand{\listingpliku}[2][]{%
  \par\smallskip\noindent
  \begin{minipage}{\linewidth}
    \lstinputlisting[#1]{#2}
  \end{minipage}
  \par\smallskip
}

\frenchspacing
\setlength{\parindent}{0.6cm}
\setlength{\parskip}{3pt}
\linespread{1.3}
\raggedbottom
\setcounter{tocdepth}{2}
\setcounter{secnumdepth}{3}
\renewcommand{\arraystretch}{1.18}
\captionsetup[table]{position=top,skip=6pt}
\captionsetup[figure]{skip=7pt}
\setlength{\emergencystretch}{2em}

\renewcommand{\cfttoctitlefont}{\Large\bfseries\sffamily}
\renewcommand{\cftloftitlefont}{\Large\bfseries\sffamily}
\renewcommand{\cftlottitlefont}{\Large\bfseries\sffamily}
\renewcommand{\cftsecfont}{\bfseries\sffamily}
\renewcommand{\cftsubsecfont}{\sffamily}

\pagestyle{fancy}
\fancyhf{}
\fancyhead[C]{\footnotesize\textit{AKADEMIA NAUK STOSOWANYCH W NOWYM SĄCZU}}
\fancyfoot[C]{%
  \scriptsize DOKUMENTACJA PROJEKTU - ZAAWANSOWANE PROGRAMOWANIE\\
  \normalsize\sffamily\thepage
}
\setlength{\headheight}{20pt}
\setlength{\footskip}{35pt}
\renewcommand{\headrulewidth}{0.4pt}
\renewcommand{\footrulewidth}{0.4pt}
\makeatletter
\let\ps@plain\ps@fancy
\makeatother
